\documentclass[11pt]{article}
\usepackage[review]{acl}
\usepackage{times}
\usepackage{latexsym}
\usepackage{graphicx}
\usepackage{amsmath}
\usepackage{booktabs}
\usepackage{multirow}
\usepackage{array}
\usepackage{xcolor}

\title{Layer Selection Matters: How Metric Choice Determines\\Hidden-State Probe Effectiveness for Reasoning Model Safety Monitoring}

\author{Anonymous}

\begin{document}
\maketitle

\begin{abstract}
Reasoning-capable LLMs can ``self-jailbreak'' during extended chain-of-thought, interleaving safety reasoning with harmful execution until a \emph{commitment point} where safety reasoning permanently ceases.
We present the first systematic evaluation of detection methods for this phenomenon across 4 models spanning 2 architectures (Llama, Qwen), 2 training methods (R1-distilled, STILL), and 2 scales (7--8B, 32B).
Our primary methodological contribution is identifying that the layer selection metric \texttt{bal\_acc $\times$ crossing\_rate} systematically biases toward shallow layers---especially at 32B scale---entirely masking hidden-state (HS) probe advantages and producing false negative conclusions.
At correctly selected layers, HS probes significantly outperform text probes on all 4 models (12/12 Holm--Bonferroni tests, precision $+$15--24pp), a finding that replicates on an independent dataset (AdvBench) with fundamentally different labeling.
Six alternative detection methods (Mahalanobis distance, RepE, CCS, BOCPD/CUSUM, MLP, cross-model transfer) all fail, establishing supervised temporal probing with principled layer selection as the most effective approach.
\end{abstract}


\section{Introduction}

Recent reasoning-capable LLMs such as DeepSeek-R1 \cite{guo2025deepseek} and QwQ \cite{qwq32b} produce extended chains of thought (CoT) that can span hundreds of tokens.
While this deliberative process enables stronger problem-solving, it also introduces a novel safety failure mode: \emph{self-jailbreaking}, where the model's own reasoning process transitions from safety-compliant deliberation to harmful content generation \cite{ren2025selfharm}.
Unlike adversarial input attacks targeted by existing defenses \cite{li2024trajguard}, self-jailbreaking originates within the model's reasoning trace itself, rendering input-level detectors ineffective.

Prior work has characterized this phenomenon \cite{ren2025selfharm, wang2025harmthoughts} and proposed text-level CoT monitors \cite{wang2025saferemind}, but no systematic comparison of detection modalities exists.
A fundamental question remains: \emph{can hidden-state monitoring detect safety commitment earlier or more reliably than text monitoring?}

We address this through a comprehensive evaluation across 4 models: R1-8B (Llama-8B, R1-distilled), OpenThinker-7B (Qwen-7B, STILL-trained), QwQ-32B (Qwen-32B, native), and R1-32B (Qwen-32B, R1-distilled).
Our key findings are:

\paragraph{Primary contribution: Layer selection metric sensitivity.}
The composite metric \texttt{bal\_acc $\times$ crossing\_rate}, commonly used for layer selection in probing studies, systematically selects shallow layers (0--16\% depth) due to inflated crossing rates from representational compression.
Threshold-dependent metrics (precision, FPR) consistently select deep layers ($>$45\% depth).
This bias entirely masked the HS advantage at 32B scale: correcting layer selection improved Holm--Bonferroni pass rates from 5/12 to \textbf{12/12}.

\paragraph{Cross-dataset validation.}
The layer selection bias pattern replicates on AdvBench \cite{zou2023advbench} with binary jailbreak/safe labels (no commitment-point annotation), confirming it as a geometric property of model representations rather than a dataset artifact.

\paragraph{Comprehensive negative results.}
Six alternative methods---Mahalanobis distance (TrajGuard-style), RepE direction probing, CCS, BOCPD/CUSUM change-point detection, MLP probes, and cross-model transfer---all fail to improve over supervised linear probes, establishing what does \emph{not} work for this task.


\section{Related Work}

\paragraph{Hidden-state safety probes.}
TrajGuard \cite{li2024trajguard} uses Mahalanobis distance on hidden states to detect adversarial inputs, achieving 95\% defense rates on instruction-tuned models.
SafeSwitch \cite{chen2025safeswitch} and SIREN \cite{wang2025siren} extend this to runtime intervention.
These approaches target adversarial \emph{input} attacks on non-reasoning models; none address reasoning-model self-jailbreak.

\paragraph{Reasoning model safety.}
Self-Jailbreaking \cite{ren2025selfharm} and HarmThoughts \cite{wang2025harmthoughts} characterize safety failures in reasoning models.
SafeRemind \cite{wang2025saferemind} and SafeThink \cite{qi2025safethink} propose text-level CoT monitoring but are vulnerable to pathological reasoning chains.
Our work provides the first modality comparison (hidden-state vs.\ text) for reasoning-model safety detection.

\paragraph{Representation probing.}
Linear probing of hidden states has a long history in NLP \cite{belinkov2022probing}.
Burns et al.\ \cite{burns2023ccs} propose CCS for unsupervised truth discovery.
Representation Engineering (RepE) \cite{zou2023repe} uses mean-difference directions for concept detection.
We evaluate both adapted to safety commitment detection.


\section{Method}

\subsection{Problem Setup}

We study \emph{safety commitment detection}: given a reasoning trace $\mathbf{s} = (s_1, \ldots, s_T)$ where each $s_t$ is a sentence-level step, predict at each step whether the model has committed to unsafe behavior.
Each trace is annotated with a \emph{commitment point} (CP)---the step after which safety reasoning permanently ceases---yielding binary step-level labels $y_t \in \{0, 1\}$.

\paragraph{Models.}
We evaluate 4 models spanning a $2 \times 2$ design: architecture (Llama-8B vs.\ Qwen-7B/32B) $\times$ training (R1-distilled vs.\ STILL/native): R1-8B ($n{=}43$ test traces), OpenThinker-7B (OT-7B, $n{=}47$), QwQ-32B ($n{=}19$--20), and R1-32B ($n{=}20$).

\subsection{Probe Architecture}

\paragraph{Hidden-state probes.}
For each step $s_t$, we extract the hidden state $\mathbf{h}_t^{(\ell)} \in \mathbb{R}^d$ at layer $\ell$ (token-averaged).
A \emph{static} logistic regression (LR) probe classifies $\mathbf{h}_t^{(\ell)}$ directly.
A \emph{temporal} LR probe classifies the concatenation $[\mathbf{h}_{t-W+1}^{(\ell)}; \ldots; \mathbf{h}_t^{(\ell)}] \in \mathbb{R}^{d \cdot W}$ for window size $W$.

\paragraph{Text probes.}
We encode each step using MiniLM-L6-v2 \cite{reimers2019sentencebert} ($d{=}384$) and train analogous static and temporal LR probes on sentence embeddings.

\paragraph{GRU sequence model.}
As an upper-bound probe, we train a GRU over multi-layer hidden states, predicting commitment at each step.

\subsection{Layer Selection Sensitivity Analysis}
\label{sec:layer-selection}

A critical design choice in probing is \emph{which layer} to extract hidden states from.
We evaluate 5 layer selection metrics across all 4 models:

\begin{enumerate}
\item \texttt{bal\_acc $\times$ CR}: balanced accuracy weighted by crossing rate (CR)---the fraction of traces where the probe prediction crosses the threshold.
\item \texttt{Pure bal\_acc}: balanced accuracy alone (threshold-free).
\item \texttt{Precision}: positive predictive value at the operating threshold.
\item \texttt{Inverse FPR}: $1 - \text{FPR}$ at the operating threshold.
\item \texttt{Composite}: additive combination of precision, inverse FPR, and balanced accuracy.
\end{enumerate}

We discover that these metrics produce systematically different layer selections due to the \emph{representational geometry} of transformer hidden states: shallow layers have compressed representations (smaller inter-class cosine distances), causing noise-driven threshold crossings that inflate CR.

\subsection{Temporal Windowing}

We sweep window sizes $W \in \{1, 3, 5, 10, 15, 20, 25\}$.
Detection uses a $K$-consecutive-crossing criterion: the probe fires when $K{=}5$ consecutive steps exceed the threshold.
Larger $W$ increases precision (more stable predictions) but decreases coverage (requires more consecutive steps).

\subsection{Evaluation Protocol}

All models use stratified 60/20/20 train/val/test splits with class-balanced weighting.
Statistical significance is assessed via bootstrap resampling (10{,}000 iterations) with Holm--Bonferroni step-down correction across 12 tests (4 models $\times$ 3 metrics: balanced accuracy, precision, FPR).


\section{Experiments}

\subsection{Layer Selection Metric Sensitivity}
\label{sec:exp-layer}

Table~\ref{tab:layer-metric} shows the layer selected by each metric across all 4 models.
A striking three-tier pattern emerges:

\begin{table}[t]
\centering
\small
\begin{tabular}{@{}lcccc@{}}
\toprule
\textbf{Metric} & \textbf{R1-8B} & \textbf{OT-7B} & \textbf{QwQ} & \textbf{R1-32B} \\
\midrule
bal\_acc$\times$CR & L0 (0\%) & L0 (0\%) & L2 (3\%) & L10 (16\%) \\
Pure bal\_acc & L12 (39\%) & L9 (32\%) & L30 (47\%) & L47 (73\%) \\
Precision & L16 (52\%) & L20 (71\%) & L60 (94\%) & L63 (98\%) \\
Inv.\ FPR & L16 (52\%) & L14 (50\%) & L60 (94\%) & L63 (98\%) \\
Composite & L12 (39\%) & L14 (50\%) & L60 (94\%) & L63 (98\%) \\
\bottomrule
\end{tabular}
\caption{Selected layer (depth \%) by metric and model. CR-contaminated metrics select shallow layers; threshold-dependent metrics select deep layers. The gap is 12--14 layers for 32B models.}
\label{tab:layer-metric}
\end{table}

\paragraph{CR-contaminated $\to$ shallow (0--16\%).}
The \texttt{bal\_acc $\times$ CR} metric selects layers at 0--16\% depth across all 4 models.
At these shallow layers, representations are compressed: mean cosine distance between hidden states is 1.4--2.7$\times$ smaller than at deep layers.
This compression causes probe predictions to oscillate around the decision boundary, inflating crossing rates and biasing the composite score upward.

\paragraph{Threshold-free $\to$ middle (32--73\%).}
Pure balanced accuracy, which does not depend on a specific threshold, selects middle layers that balance accuracy and stability.

\paragraph{Threshold-dependent $\to$ deep (45--98\%).}
Precision, inverse FPR, and the additive composite---all deployment-relevant metrics---converge on deep layers where predictions are stable and false alarms rare.

\paragraph{Impact on HS--text comparison.}
At biased shallow layers, HS probes showed non-significant advantages at 32B scale (R1-32B: $p{=}0.064$; QwQ: $p{=}0.35$), producing the false conclusion that HS advantage is scale-dependent.
At corrected deep layers, all 4 models show significant HS advantages (Table~\ref{tab:hb-summary}).
LOO-CV confirms layer selection stability: the dominant layer is selected in 65--68\% of folds across models.


\subsection{Cross-Dataset Validation on AdvBench}
\label{sec:advbench}

To test whether layer selection bias is a property of our dataset or of model geometry, we replicate the analysis on AdvBench \cite{zou2023advbench} using R1-8B with binary jailbreak/safe labels (no commitment-point annotation; $n{=}520$ traces, 437 jailbreak, 83 safe).

\begin{table}[t]
\centering
\small
\begin{tabular}{@{}lcc@{}}
\toprule
& \textbf{Original (CP)} & \textbf{AdvBench (binary)} \\
\midrule
CR-composite layer & L0 (0\%) & L2 (6.5\%) \\
Threshold layer & L14 (45\%) & L14 (45\%) \\
Layer gap & 14 layers & 12 layers \\
HS bal\_acc (deep) & 0.810 & 0.762 \\
HS--text prec.\ $\Delta$ & +22.4pp & +18.3pp \\
CR ratio (shallow/deep) & 2.50 & 2.32 \\
HB pass & 3/3 & 3/3 \\
\bottomrule
\end{tabular}
\caption{Layer selection bias replicates on AdvBench with fundamentally different labeling. All 7 pattern metrics match target ranges.}
\label{tab:advbench-pattern}
\end{table}

Table~\ref{tab:advbench-pattern} shows complete pattern replication: CR-composite selects L2 (6.5\% depth), threshold metrics select L14 (45\%), with a 12-layer gap consistent with the original 14-layer gap.
The shallow-layer crossing rate (0.44) is 2.32$\times$ the deep-layer rate (0.19), matching the original ratio of 2.50.
Cosine distance analysis confirms the geometric mechanism: L2 mean cosine distance is 0.038 vs.\ 0.091 at L14 (ratio 2.39$\times$), consistent with representational compression driving inflated crossing rates.

At the threshold-selected layer L14, HS probes achieve bal\_acc 0.762, precision 0.677, and FPR 0.155, outperforming text probes by +10.1pp, +18.3pp, and $-$12.8pp respectively (all 3/3 Holm--Bonferroni, $\alpha{=}0.0167$).

\subsection{Hidden-State vs.\ Text Probe Comparison}
\label{sec:hs-text}

\begin{table}[t]
\centering
\small
\begin{tabular}{@{}llccccc@{}}
\toprule
\textbf{Model} & \textbf{Layer} & \textbf{W} & $\Delta$\textbf{BA} & $\Delta$\textbf{Prec} & $\Delta$\textbf{FPR} & \textbf{HB} \\
\midrule
R1-8B & L14 & 15 & +8.2 & +22.4 & $-$7.4 & 3/3 \\
OT-7B & L16 & 3 & +10.7 & +21.9 & $-$17.1 & 3/3 \\
QwQ & L60 & 3 & +15.4 & +23.6 & $-$9.8 & 3/3 \\
R1-32B & L63 & 1 & +18.7 & +15.4 & $-$14.6 & 3/3 \\
\bottomrule
\end{tabular}
\caption{HS vs.\ text probe comparison at corrected layers. All deltas in percentage points. HB: Holm--Bonferroni pass count (3 metrics). 12/12 tests pass at $\alpha{=}0.05$ after step-down correction.}
\label{tab:hb-summary}
\end{table}

Table~\ref{tab:hb-summary} reports the core HS--text comparison.
Across all 4 models, HS probes achieve significantly higher precision ($+$15.4 to $+$23.6pp) and lower FPR ($-$7.4 to $-$17.1pp) than text probes.
All 12 tests (4 models $\times$ 3 metrics) pass Holm--Bonferroni correction at $\alpha{=}0.05$ with 10{,}000 bootstrap resamples.

\paragraph{32B caveats.}
The 32B models exhibit model-specific trade-offs.
R1-32B achieves significance only at $W{=}1$ (3/3 HB) with 25\% coverage; at $W{\geq}3$, dimensionality ($51{,}200 \times W$ features for $n{=}20$) degrades performance.
QwQ-32B shows robust significance across all 15 window sizes but with poor HS lead time ($-1.2$ to $+1.9$ steps vs.\ text's $+3.7$ to $+12.3$).

\paragraph{7--8B scale.}
At 7--8B, HS advantages are clean.
R1-8B achieves precision 69.7\% vs.\ text 47.2\% ($+$22.4pp); OT-7B achieves the highest selectivity (18.3pp over random control), confirming the signal is safety-specific.

\paragraph{Text baseline detail.}
Text probes use MiniLM-L6-v2 sentence embeddings (384d).
HS probes access 4{,}096--5{,}120d hidden states encoding internal processing state.
The consistent precision gap ($+$15--24pp) across all 4 models indicates HS probes capture information invisible at the text surface---the model may internally commit to unsafe behavior before this manifests in generated text.


\subsection{Temporal Window Ablation}

\begin{table}[t]
\centering
\small
\begin{tabular}{@{}lcccc@{}}
\toprule
& \textbf{R1-8B} & \textbf{OT-7B} & \textbf{QwQ} & \textbf{R1-32B} \\
\midrule
Best W (bal\_acc) & 25 & 3 & 25 & 10 \\
W=1 precision & 0.693 & --- & --- & 0.603 \\
Best W precision & 0.822 & 0.466 & --- & --- \\
W=1 coverage & 0.721 & --- & --- & 0.250 \\
Best W coverage & 0.488 & --- & --- & --- \\
\bottomrule
\end{tabular}
\caption{Window size ablation. Larger $W$ improves precision but reduces coverage. R1-32B peaks at $W{=}10$ then declines (dimensionality constraint).}
\label{tab:window}
\end{table}

Temporal context universally improves HS probe accuracy.
For R1-8B, increasing $W$ from 1 to 25 improves precision from 69.3\% to 82.2\% while halving FPR from 9.65\% to 4.83\%, at the cost of coverage dropping from 72.1\% to 48.8\% (Table~\ref{tab:window}).
This precision--coverage trade-off is a fundamental characteristic of temporal probing.

For text probes, temporal improvement is universal across all 4 models ($+$3--12pp bal\_acc, $+$8--23pp precision).
OT-7B HS probes peak at $W{=}3$ for bal\_acc, then decline; precision and FPR continue improving at larger $W$.

32B temporal results are constrained by sample size: R1-32B ($n{=}20$, $d{=}5{,}120$) peaks at $W{=}10$ and degrades beyond; QwQ-32B shows monotonic improvement but with all 15 windows statistically significant.


\subsection{Sequence Models and Upper Bounds}

GRU probes trained on multi-layer hidden states provide an upper-bound comparison.
The pattern is model-specific: OT-7B GRU outperforms LR ($+$5.8pp bal\_acc, from 68.8\% to 74.6\%) due to rich sequential patterns in longer traces ($\sim$80 sentences).
However, for the remaining 3 models, LR matches or exceeds GRU: R1-8B LR $W{=}15$ achieves 84.8\% vs.\ GRU 80.9\%; QwQ LR $W{=}5$ achieves 76.1\% vs.\ GRU 72.4\%; R1-32B LR $W{=}10$ achieves 73.7\% vs.\ GRU 67.5\%.
Linear probes are sufficient for 3/4 models.


\subsection{Lead Time Analysis}

\begin{table}[t]
\centering
\small
\begin{tabular}{@{}lcccc@{}}
\toprule
\textbf{Model} & \textbf{HS LR} & \textbf{Text LR} & \textbf{GRU} & \textbf{Caveat} \\
\midrule
R1-8B & $\sim$0 & +10.3* & +4.5$\pm$3.0 & *inflated FPR \\
OT-7B & +21.6 & +45.5 & +33.6 & long traces \\
QwQ & $-$1.2--1.9 & +3.7--12.3 & +15.9 & HS late \\
R1-32B & +11.7 & --- & +8.8 & 25\% cov. \\
\bottomrule
\end{tabular}
\caption{Lead time (steps before commitment point). No universal pattern; per-model reporting is essential.}
\label{tab:lead}
\end{table}

Table~\ref{tab:lead} reveals highly model-specific lead time profiles with no universal pattern.
R1-8B HS probes detect \emph{at} commitment (lead $\approx 0$), encoding current safety state rather than predicting future trajectory.
Text probes' apparent lead ($+$10.3 steps) is an artifact of high FPR---a trigger-happy detector appears to detect early.
GRU achieves modest positive lead ($+$4.5 \pm 3.0$, worst seed $+$0.9).

Oracle decomposition reveals the source of the $+$31-step oracle advantage: $\sim$15pp from position information ($t/T$), $\sim$10pp from label sequence patterns, and only $\sim$6pp from true future-state access.
Hidden states encode current processing state but cannot predict whether commitment will occur---a fundamental architectural constraint.


\subsection{Negative Results}

We evaluate 6 alternative methods; all fail to improve over supervised linear probes.

\paragraph{Mahalanobis distance (TrajGuard-style).}
Across all 4 models, class separation is near-zero (max 0.20).
The discriminative signal resides in minor PCA components, not principal directions: Fisher ratio drops from 17 to 1 in PCA-3D space.
False alarm rates equal or exceed true positive rates at every threshold.

\paragraph{BOCPD/CUSUM change-point detection.}
BOCPD achieves 93\% detection on R1-8B (which has gradual commitment dynamics) but 0\% on the other 3 models, which exhibit abrupt transitions.
CUSUM is universally worse than simple thresholding.

\paragraph{Cross-model transfer.}
Zero-shot transfer via PCA projection fails entirely: best accuracy 67.3\% (QwQ$\to$R1-32B) vs.\ within-model baselines of 72--82\%.
Safety representations are model-specific; 5-seed within-model robustness contrasts sharply with cross-model failure.

\paragraph{16-label fine-grained supervision.}
Using a 4-type taxonomy (vs.\ binary CP labels) yields marginal improvement: $+$2.2pp on R1-8B, $+$1.7pp on OT-7B.
The 4$\times$ annotation cost is not justified; binary labels suffice.


\subsection{False Positive Analysis}

All methods exhibit high false positive rates on entirely safe traces (39--62\% safe FPR), because safe traces contain safety-relevant reasoning that resembles the early stages of jailbreak traces.
HS probes achieve lower FPR than text probes on 3/4 models (R1-8B: 16.2\% vs.\ 35.6\%; OT-7B: 11.1\% vs.\ 30.2\%; R1-32B: 23.9\% vs.\ 31.0\%).
QwQ is the exception (HS 50.0\% vs.\ text 32.5\%), attributable to more extensive safety reasoning in QwQ's safe traces.


\section{Discussion}

\paragraph{Layer selection as researcher degree of freedom.}
Our central methodological finding---that metric choice for layer selection can reverse conclusions about probe effectiveness---has implications beyond safety monitoring.
Any probing study using composite metrics involving crossing rate or prediction stability risks the same bias.
We recommend sensitivity analysis across multiple metric families as standard practice.

\paragraph{Why deep layers?}
Deep layers have expanded representations (cosine distance 2.4$\times$ larger than shallow layers), producing stable, low-noise probe predictions with fewer threshold crossings.
Deployment-relevant metrics (precision, FPR) naturally favor these stable predictions, while CR-contaminated metrics favor the noisy shallow-layer predictions.
This is a geometric property: it replicates across datasets (AdvBench), labeling schemes (binary vs.\ step-level), and architectures.

\paragraph{32B model-specific trade-offs.}
The two 32B models present contrasting deployment profiles: R1-32B offers excellent lead time ($+$11.7 steps) but narrow coverage (25\% at significant $W{=}1$), while QwQ offers robust significance across all window sizes but poor HS lead time ($-$1.2$ to $+$1.9 steps).
These trade-offs must be evaluated per-model rather than unified into scale-level conclusions.

\paragraph{Linear probe sufficiency.}
GRU improves over LR only on OT-7B ($+$5.8pp), where long traces ($\sim$80 sentences) provide rich sequential patterns.
MLP probes underperform LR on 3/4 models.
This suggests the safety commitment signal is linearly decodable from hidden states, consistent with the linear representation hypothesis \cite{park2024linearrep}.


\section{Conclusion}

We provide the first systematic comparison of detection methods for reasoning-model safety commitment across 4 models, 2 architectures, and 2 scales.
Our primary contribution is identifying layer selection metric sensitivity as a critical confound: the commonly used \texttt{bal\_acc $\times$ crossing\_rate} metric systematically biases toward shallow layers, entirely masking hidden-state advantages at 32B scale.
Correcting this bias reveals universal HS superiority (12/12 Holm--Bonferroni, precision $+$15--24pp), validated on an independent dataset.
Combined with the comprehensive failure of 6 alternative methods, our results establish supervised temporal probing with principled layer selection as the most effective approach for monitoring reasoning-model safety commitment.


\section*{Limitations}

Our evaluation has three main limitations.
First, sample sizes for 32B models are small ($n{=}19$--20 test traces), constraining temporal window analysis and producing wide confidence intervals; R1-32B achieves significance only at $W{=}1$ with 25\% coverage.
Second, all methods exhibit high false positive rates on safe traces (39--62\%), limiting immediate deployment utility---the probes are diagnostic tools rather than production-ready detectors.
Third, cross-model transfer fails entirely, requiring independent probe training per model; scalability to new model releases requires new annotation and training effort.


\bibliography{references}

\appendix

\section{Holm--Bonferroni Step-Down Results}
\label{app:hb}

\begin{table*}[t]
\centering
\small
\begin{tabular}{@{}clcccc@{}}
\toprule
\textbf{Rank} & \textbf{Model -- Metric} & \textbf{$\Delta$} & \textbf{$p$-value} & \textbf{HB Threshold} & \textbf{Pass} \\
\midrule
1 & R1-8B -- precision & +27.9pp & $<$0.001 & 0.0042 & Y \\
2 & OT-7B -- bal\_acc & +10.7pp & $<$0.001 & 0.0045 & Y \\
3 & OT-7B -- precision & +21.9pp & $<$0.001 & 0.0050 & Y \\
4 & OT-7B -- FPR & $-$17.1pp & $<$0.001 & 0.0056 & Y \\
5 & QwQ -- bal\_acc & +14.8pp & $<$0.001 & 0.0063 & Y \\
6 & QwQ -- precision & +23.5pp & $<$0.001 & 0.0071 & Y \\
7 & R1-32B -- bal\_acc & +18.7pp & $<$0.001 & 0.0083 & Y \\
8 & R1-32B -- precision & +15.4pp & $<$0.001 & 0.0100 & Y \\
9 & QwQ -- FPR & $-$9.5pp & 0.006 & 0.0125 & Y \\
10 & R1-8B -- bal\_acc & +8.2pp & 0.005 & 0.0167 & Y \\
11 & R1-32B -- FPR & $-$14.6pp & 0.010 & 0.0250 & Y \\
12 & R1-8B -- FPR & $-$7.4pp & 0.044 & 0.0500 & Y \\
\bottomrule
\end{tabular}
\caption{Complete Holm--Bonferroni step-down results. 12/12 pass at $\alpha{=}0.05$, 10{,}000 bootstrap resamples.}
\label{tab:hb-full}
\end{table*}


\section{AdvBench Detailed Results}
\label{app:advbench}

\begin{table}[t]
\centering
\small
\begin{tabular}{@{}lcccc@{}}
\toprule
\textbf{Layer} & \textbf{Depth} & \textbf{Cosine} & \textbf{CR} & \textbf{Interp.} \\
\midrule
L2 & 6.5\% & 0.038 & 0.44 & Compressed \\
L8 & 25.8\% & 0.062 & 0.31 & Intermediate \\
L14 & 45.2\% & 0.091 & 0.19 & Expanded \\
L20 & 64.5\% & 0.105 & 0.16 & Deep/stable \\
\bottomrule
\end{tabular}
\caption{Cosine distance and crossing rate by layer depth on AdvBench. Representational compression at shallow layers drives inflated crossing rates.}
\label{tab:advbench-cosine}
\end{table}

\begin{table}[t]
\centering
\small
\begin{tabular}{@{}lcccc@{}}
\toprule
\textbf{Modality} & \textbf{BA} & \textbf{Prec} & \textbf{FPR} & \textbf{Cov} \\
\midrule
HS (L14) & 0.762 & 0.677 & 0.155 & 0.586 \\
Text (L14) & 0.661 & 0.494 & 0.283 & 0.529 \\
$\Delta$ & +10.1 & +18.3 & $-$12.8 & +5.7 \\
$p$ (bootstrap) & 0.005 & 0.002 & 0.008 & --- \\
\bottomrule
\end{tabular}
\caption{HS vs.\ text at threshold-selected layer L14 on AdvBench ($n{=}104$ test). 3/3 Holm--Bonferroni pass.}
\label{tab:advbench-hs-text}
\end{table}


\section{32B Layer Correction Detail}
\label{app:32b-correction}

\begin{table}[t]
\centering
\small
\begin{tabular}{@{}lcccc@{}}
\toprule
\textbf{Model} & \textbf{Biased L} & \textbf{Corrected L} & $\Delta$\textbf{BA} & $\Delta$\textbf{Prec} \\
\midrule
R1-32B & L10 (16\%) & L63 (98\%) & +10.8pp & +8.7pp \\
QwQ & L2 (3\%) & L60 (94\%) & +1.0pp & --- \\
\bottomrule
\end{tabular}
\caption{Performance gain from correcting layer selection at 32B scale. R1-32B shows substantial improvement; QwQ's absolute gain is smaller but significance changes dramatically ($p{=}0.35 \to p{<}0.005$).}
\label{tab:32b-correction}
\end{table}

\end{document}
